\section{A Guided Explanation and Notation}
\label{app:tutorial}
\subsection{The problem in plain language}
Imagine answering a question while looking at a long collection of notes.
Keeping every note makes it easy to look things up but fills the desk.
Throwing away the least-used notes saves space, yet a rarely used note might
contain the one number needed for the final answer. MosaicKV is a way of
choosing which internal notes a vision--language model keeps on its GPU.
The analogy has a limit: the cache does not contain human-readable facts.
It contains learned numerical vectors whose importance can be difficult to
interpret.

A \emph{token} is one model input position. Text can be split into pieces of
words. An image is converted into visual features that contribute many input
positions. A \emph{key} helps decide whether a position is relevant to the
current request; a \emph{value} is the information mixed into the attention
output if that position is attended to. A \emph{query} is the vector asking
what information is useful now. These descriptions are intuitions, not literal
symbols stored in memory.

During \emph{prefill}, the model reads the image and prompt. During
\emph{decoding}, it predicts one output token at a time. Caching saves the
model from re-reading the entire prefix at every step. Compression inserts a
decision between these stages: keep a smaller useful state while preserving
the model's ability to answer.

\subsection{Why each component exists}
\paragraph{Forecasting: anticipate the next question.}
Prompt attention says what mattered while reading the input. Future-query
forecasting tries to estimate what will matter while producing the answer.
A draft is like briefly rehearsing an answer without committing it. Its
cost is real and must be included in runtime.

\paragraph{Graph: keep related pieces together.}
A document answer may need both a row label and its number. Selecting each
block independently can preserve one and discard the other. The graph makes
relationships available to selection. It cannot recover OCR text or geometry
that the model adapter never exposed.

\paragraph{Value awareness: importance is more than attention.}
Two blocks might receive the same attention but produce very different value
vectors. Conversely, two frequently attended blocks may carry nearly duplicate
values. A value-aware score tries to distinguish contribution from repetition.
It is still a surrogate for final answer correctness.

\paragraph{Budgeted selection: make the decision fit.}
The selector must solve a constrained choice, not merely rank items and hope
the result fits. The budget includes physical allocation effects. Mandatory
tokens and protected regions consume capacity before optional choices.

\paragraph{Tiers: distinguish preservation, approximation, and recovery.}
Exact storage keeps originals. Prototypes summarize compatible discarded
members. Residual storage retains originals for potential recovery. A summary
can be wrong; an original on the CPU still occupies memory and needs time to
transfer. None of these costs disappears because a tier has a different name.

\paragraph{Repair: spend extra work only when useful.}
Repair attempts to recover evidence when a provisional step appears uncertain.
An uncertainty signal is not an oracle. A confident wrong answer may never
trigger, and an uncertain correct answer may cause unnecessary restoration.
Both false positives and missed beneficial repairs belong in the evaluation.

\subsection{Glossary}
\begin{table}[h]
\caption{Abbreviations used throughout the manuscript.}
\small\begin{tabularx}{\columnwidth}{lY}\toprule
Term & Meaning\\\midrule
KV & Key--value; stored attention keys and values\\
VLM / LLM & Vision--language / large language model\\
GPU / CPU & Graphics / central processing unit\\
HBM & High-bandwidth memory on the accelerator\\
HF & Hugging Face model/runtime ecosystem\\
GQA & Grouped-query attention\\
MHA & Multi-head attention\\
RoPE & Rotary position embedding\\
SDPA & Scaled dot-product attention\\
OCR & Optical character recognition\\
CSR / COO & Compressed sparse row / coordinate storage\\
FP32 / FP16 & 32-bit / 16-bit floating point\\
BF16 & Brain floating point, a 16-bit format\\
TTFT & Time to first generated token\\
ITL & Inter-token latency\\
KL & Kullback--Leibler divergence\\
ANLS & Average normalized Levenshtein similarity\\
PPL / NLL & Perplexity / negative log-likelihood\\
CI & Confidence interval\\
MiB / GiB & $2^{20}$ / $2^{30}$ bytes\\
SHA & Secure hash algorithm; artifact/source identifier\\
JSONL & One JSON record per line\\
\bottomrule\end{tabularx}
\end{table}

\subsection{Notation}
\begin{table}[h]
\caption{Core mathematical symbols.}
\small\begin{tabularx}{\columnwidth}{lY}\toprule
Symbol & Definition\\\midrule
$\ell,h$ & Layer and KV-head index\\
$N,n$ & Source token count and number of blocks\\
$b,d,s$ & Batch size, head dimension, element bytes\\
$z$ & Configured token-block size\\
$\mathcal B,\mathcal M$ & Source blocks and mandatory subset\\
$S,C$ & Selected set and per-sample byte cap\\
$w,r$ & Prompt-window length and draft steps\\
$a_i,c_i^v,n_i$ & Attention, value contribution, novelty\\
$w_{ij}$ & Sparse nonnegative graph edge weight\\
$P,B_{\mathrm{active}}(P)$ & Cache plan and its physical active bytes\\
$p_j,\pi(j)$ & Logical position and physical-to-logical map\\
$R$ & Maximum residual blocks restored per event\\
\bottomrule\end{tabularx}
\end{table}

\subsection{A position example, not an experimental result}
Figure~\ref{fig:logical} shows eight illustrative source positions and four
retained entries. This is a schematic, not a measurement. After compaction,
physical slot one contains original position two. The next generated token
still uses logical position eight. Renumbering retained entries as positions
zero through three would change their positional meaning.

\begin{figure}[h]
\centering\begin{tikzpicture}[x=0.83cm,y=0.85cm,font=\small]
\node[anchor=west] at (-0.45,2.8) {\textbf{Source cache: original positions}};
\foreach \p in {0,1,2,3,4,5,6,7}{
 \node[draw,minimum width=0.7cm,minimum height=0.55cm,fill=gray!8] at (\p,2) {\p};}
\node[anchor=west] at (-0.45,-0.65) {\textbf{Packed exact cache}};
\foreach \x/\p in {0/0,1/2,2/5,3/7}{
 \node[draw=exactblue,fill=exactblue!10,minimum width=0.7cm,minimum height=0.55cm] at (\x,0) {\p};
 \draw[flow] (\p,1.62)--(\x,0.38);}
\node[align=left,anchor=west,font=\scriptsize] at (4.0,0)
 {active length: 4\\logical length: 8\\next position: 8};
\node[font=\scriptsize,anchor=west] at (-0.45,-1.2)
 {Stored labels are logical positions, not new RoPE positions.};
\end{tikzpicture}

\caption{Illustrative physical compaction with original positions preserved.
The numeric labels describe a toy layout only.}
\label{fig:logical}
\end{figure}

\subsection{What training-free does and does not mean}
The compression path does not optimize model weights or train an adapter.
It may select fixed hyperparameters on a disjoint calibration set. Such
selection must be disclosed and frozen before evaluation. No task-label fitting
or gradient updates are introduced under the MosaicKV name. Access to a
FullKV reference for a bounded online diagnostic must also be disclosed: it
adds inference work even though it is not training.
